\chapter{Trigger finger release}
\index{Trigger finger release}

\section*{Definition and Overview}
Trigger finger, known in the clinical literature as stenosing tenosynovitis, is a common and often painful hand pathology characterised by the impairment of normal tendon gliding within the digital flexor sheath. The human hand relies on a complex apparatus of muscles, tendons, and ligaments to perform precise movements. The flexor tendons---specifically the flexor digitorum superficialis (FDS) and the flexor digitorum profundus (FDP)---originate in the forearm and travel through the carpal tunnel into the palm, inserting into the phalanges of the fingers. To prevent these tendons from pulling away from the bone during flexion (a phenomenon known as bowstringing), they are held securely against the phalanges by a specialised retinacular tunnel known as the flexor tendon sheath. This sheath is composed of a series of thickened annular (A1 to A5) and cruciate (C1 to C3) pulleys that act as mechanical pulleys, optimising the excursion and force transmission of the tendons.

The pathophysiology of trigger finger focuses on the A1 pulley, which lies directly volar to the metacarpophalangeal (MCP) joint. Due to localised mechanical stress, repetitive microtrauma, or systemic metabolic changes, a mismatch develops between the volume of the flexor tendons and the diameter of the A1 pulley tunnel. This mismatch manifests as localised swelling, sheath thickening, or the formation of a nodular expansion within the tendon itself. When the patient attempts to flex the digit, the swollen tendon or nodule passes through the A1 pulley into a proximal position. However, when attempting to extend the finger, the nodule becomes mechanically blocked at the proximal entrance of the A1 pulley. This results in a classic catching, clicking, or locking sensation, where the finger remains stuck in a flexed posture.

Trigger finger release (TFR) is the surgical procedure designed to permanently resolve this mechanical obstruction. When conservative therapies such as orthotic splinting, activity modification, non-steroidal anti-inflammatory drugs (NSAIDs), or local corticosteroid injections fail to yield long-term improvement, surgical intervention becomes necessary. The primary objective of trigger finger release is to surgically divide the A1 pulley longitudinally, thereby widening the tendon sheath and allowing the FDS and FDP tendons to glide freely without resistance. TFR is one of the most common hand operations performed globally and can be executed via an open surgical approach or a percutaneous technique under local anaesthesia. It represents a highly effective intervention that restores manual dexterity, alleviates chronic discomfort, and prevents the development of irreversible joint contractures.

\section*{Epidemiology}
Stenosing tenosynovitis is a highly prevalent condition within the general population, with a lifetime incidence estimated at approximately 2\% to 3\%. The epidemiological profile of the disease reveals specific vulnerability patterns across demographics, age groups, and patients with pre-existing medical comorbidities.

The condition exhibits a bimodal age distribution. The primary peak occurs in adults between their fifth and sixth decades of life, with a median age of diagnosis between 50 and 60 years. In this population, the condition is usually acquired and is highly associated with repetitive hand use or systemic illness. A secondary, much smaller peak is observed in early childhood, typically before the age of two. This paediatric variant, termed congenital or paediatric trigger thumb, almost exclusively affects the thumb and is characterised by a palpable nodule (Notta's node) in the flexor pollicis longus (FPL) tendon.

Gender is a significant demographic factor in the epidemiology of trigger finger. Females are affected far more frequently than males, with a reported ratio ranging from 2:1 to as high as 6:1. While the precise physiological reasons for this disparity remain under investigation, researchers suggest it may be linked to hormonal fluctuations, particularly in postmenopausal women, as well as distinct anatomical proportions of the flexor sheath and differences in occupational manual exposure.

The dominant hand is involved in the majority of cases, reflecting the role of cumulative mechanical workload in the pathogenesis of the condition. Among the digits, the ring finger and the thumb are the most frequently affected, followed by the middle finger, the index finger, and, least commonly, the little finger. Although many patients present with a single affected digit, it is not uncommon for individuals---particularly those with underlying metabolic diseases---to present with multiple involved digits, either simultaneously or sequentially. In patients diagnosed with diabetes mellitus, the lifetime risk of developing trigger finger escalates dramatically to approximately 10\% to 20\%, and these individuals are significantly more prone to multi-digit involvement and resistance to non-surgical therapies.

\section*{Aetiology and Pathophysiology}
The development of stenosing tenosynovitis is multifactorial, arising from a combination of mechanical, histological, and systemic factors. The primary pathological process is not a classic acute inflammation, but rather a degenerative, hypertrophic proliferative process of the tendon sheath. The key factors involved in its aetiology and pathophysiology include:

\begin{itemize}
  \item \textbf{Repetitive Mechanical Shear Stress}: Continuous, repetitive activities requiring forceful gripping, pinching, or the manipulation of tools (especially those that vibrate or have hard edges that press directly against the palmar creases) generate high shear forces at the A1 pulley. This chronic mechanical friction between the flexor tendons and the A1 pulley triggers a cellular stress response.
  \item \textbf{Fibrocartilaginous Metaplasia}: In response to persistent friction, the normal tissue structure of the A1 pulley undergoes a histopathological transition. The inner gliding layer of the pulley, normally composed of spindle-shaped fibroblasts and collagen fibres, undergoes chondroid or fibrocartilaginous metaplasia. This results in the accumulation of extracellular matrix proteins, glycosaminoglycans, and chondrocytes, leading to significant thickening, rigidity, and loss of compliance of the pulley wall.
  \item \textbf{Tenosynovial Hypertrophy and Nodule Formation}: The friction simultaneously affects the flexor tendons and their surrounding tenosynovial sheath. The tenosynovium becomes congested, hyperaemic, and hypertrophic. The tendon itself may develop localised mucoid degeneration, leading to diameter expansion or the formation of a distinct fibrotic nodule. This nodule (often referred to as Notta's node in children) creates a focal mismatch between the tendon and the pulley.
  \item \textbf{Diabetes Mellitus and Glycation End-Products}: The high prevalence of trigger finger in diabetic individuals is linked to systemic biochemical alterations. Elevated blood glucose levels promote the formation of advanced glycation end-products (AGEs), which form abnormal covalent cross-links between collagen fibres. This process stiffens the tendons and the A1 pulley, reducing their natural elasticity and predisposing them to mechanical impingement.
  \item \textbf{Inflammatory Arthropathies}: Systemic inflammatory conditions such as rheumatoid arthritis, psoriatic arthritis, and gout cause proliferative synovitis within the flexor tendon sheath. In rheumatoid arthritis, the synovium proliferates aggressively (pannus formation), filling the sheath and restricting tendon movement, which can cause triggering at the A1 pulley or occasionally at the more distal A2 or A3 pulleys.
  \item \textbf{Anatomical Predispositions}: Individuals with pre-existing stenosing conditions, such as carpal tunnel syndrome or De Quervain's tenosynovitis, have a higher statistical likelihood of developing trigger finger, suggesting a genetic or structural vulnerability to tendon-sheath mismatch.
  \item \textbf{Metabolic and Endocrine Disorders}: Hypothyroidism, amyloidosis, and renal failure are associated with trigger finger due to the systemic deposition of mucopolysaccharides, amyloid proteins, or fluid accumulation within the hand's fascial compartments, which compresses the flexor sheath.
\end{itemize}

\section*{Clinical Features}
The clinical presentation of trigger finger varies depending on the severity and duration of the pathology. Symptoms generally start insidiously and follow a progressive trajectory. The primary clinical features include:

\begin{itemize}
  \item \textbf{Palmar Pain and Localised Tenderness}: The earliest symptom is usually a dull, aching pain or discomfort at the base of the affected finger or thumb in the palm. Upon examination, there is exquisite tenderness to palpation directly over the A1 pulley, which lies near the distal palmar crease for the fingers and the MCP joint crease for the thumb.
  \item \textbf{Morning Stiffness and Hesitation}: Patients frequently report that the affected digit is stiff, painful, and difficult to move upon waking. This morning exacerbation occurs because the hand remains inactive during sleep, leading to localised fluid accumulation (oedema) within the tendon sheath, which temporarily worsens the mechanical mismatch. As the patient begins to use the hand, the fluid is gradually redistributed, and the stiffness partially resolves.
  \item \textbf{Audible or Palpable Clicking}: As the patient flexes and extends the digit, a distinct snapping, clicking, or popping sensation is felt and sometimes heard. This represents the thickened portion of the tendon or nodule forcing its way through the narrowed A1 pulley.
  \item \textbf{Mechanical Locking}: As the disease progresses, the tendon nodule becomes trapped proximal to the A1 pulley during active flexion. The finger becomes locked in a bent position. Initially, the patient can actively override this block with a sudden, painful snap. Later, the patient cannot actively extend the finger and must use their other hand to physically pull the digit straight.
  \item \textbf{Palpable Nodule}: A discrete, firm, and tender nodule may be palpable in the palm, immediately proximal to the MCP joint. This nodule is dynamic, moving longitudinally in tandem with the flexion and extension of the flexor tendons.
  \item \textbf{Stages of Severity (Green's Classification)}:
  \begin{itemize}
    \item \textit{Grade I (Pre-triggering)}: Pain and localised tenderness over the A1 pulley, but no objective catching or locking of the digit.
    \item \textit{Grade II (Active)}: Active triggering is present. The patient experiences demonstrable catching or clicking but can actively extend the digit without assistance.
    \item \textit{Grade III (Passive)}: Mechanical locking is present. The patient is unable to actively extend the digit and requires passive assistance from their other hand to straighten it (Grade IIIa), or is unable to actively flex the finger fully due to tendon entrapment (Grade IIIb).
    \item \textit{Grade IV (Contracture)}: Fixed flexion contracture. The proximal interphalangeal (PIP) joint has become chronically stiffened, and the digit cannot be straightened even with passive assistance from the examiner.
  \end{itemize}
\end{itemize}

\section*{Differential Diagnosis}
When evaluating a patient with finger stiffness, pain, or restricted movement, it is essential to distinguish trigger finger from other musculoskeletal, inflammatory, and neurological conditions of the hand. The differential diagnosis includes:

\begin{itemize}
  \item \textbf{Dupuytren's Contracture}: This is a progressive fibroproliferative disorder of the palmar fascia that leads to fixed flexion deformities, most commonly affecting the ring and little fingers. Unlike trigger finger, there is no mechanical catching, clicking, or dynamic locking. On palpation, Dupuytren's contracture presents as firm, non-tender nodules and thick cords within the palmar fascia that are adherent to the skin and do not move during tendon excursion.
  \item \textbf{Metacarpophalangeal (MCP) Joint Osteoarthritis}: Degenerative joint disease of the MCP joint can cause localised pain, swelling, and a mechanical block to range of motion. However, it is characterised by joint crepitus rather than snapping, and radiographs will confirm joint space narrowing, subchondral sclerosis, and osteophytes.
  \item \textbf{Flexor Tendon Rupture}: The sudden or gradual rupture of the FDS or FDP tendons results in a complete inability to actively flex the digit at the PIP or distal interphalangeal (DIP) joints. While active motion is lost, the passive range of motion remains completely free, and there is no mechanical resistance or popping sensation.
  \item \textbf{Infectious Tenosynovitis}: A rapid, bacterial infection of the flexor tendon sheath, which is a limb-threatening surgical emergency. It is clinically differentiated by the presence of Kanavel's four cardinal signs: the digit is held in a flexed posture; there is uniform, symmetrical swelling of the entire digit (often described as a ``sausage digit''); there is exquisite tenderness along the entire course of the flexor sheath; and severe pain is elicited upon passive extension of the digit.
  \item \textbf{MCP Joint Locked or Dislocated}: Mechanical locking of the joint itself can occur due to entrapment of the volar plate, collateral ligament tears, or loose intra-articular bodies. This block is constant, does not vary with tendon movement, and is often associated with a history of acute trauma.
  \item \textbf{De Quervain's Tenosynovitis}: This condition involves stenosing tenosynovitis of the first dorsal compartment of the wrist, affecting the abductor pollicis longus and extensor pollicis brevis. It is distinguished by pain localised to the radial styloid and a positive Finkelstein's test, rather than pain and locking at the base of the fingers or palmar aspect of the thumb.
  \item \textbf{Rheumatoid Arthritis and Flexor Tenosynovitis}: While rheumatoid arthritis can cause trigger finger, diffuse rheumatoid tenosynovitis without focal A1 pulley locking can cause general hand stiffness and limited tendon excursion. The presence of systemic inflammatory symptoms, symmetrical joint involvement, and laboratory markers help differentiate this condition.
\end{itemize}

\section*{When to Seek Medical Care}
Trigger finger is generally a non-emergent, chronic condition, but timely medical evaluation is essential to prevent permanent joint stiffness. Patients should schedule a consultation with a primary care physician or a hand specialist if they experience persistent pain, swelling, stiffness, or catching in their fingers that does not respond to rest or conservative measures within a few weeks.

Immediate, emergency medical care is mandatory if the patient exhibits any signs of a hand infection. Hand infections, particularly infectious flexor tenosynovitis, can spread rapidly within the confined spaces of the hand, risking permanent tendon necrosis, loss of hand function, or systemic sepsis. Emergency signs include:
\begin{itemize}
  \item Rapidly spreading redness, warmth, or red streaks travelling up the hand or arm.
  \item Severe, throbbing pain that worsens progressively and does not respond to elevation or simple analgesics.
  \item Significant, uniform swelling of the entire finger.
  \item Inability to move or straighten the finger due to intense pain.
  \item Pus or foul-smelling drainage from a wound or the palm.
  \item Systemic symptoms such as fever, chills, sweating, nausea, or confusion.
\end{itemize}

If any of these emergency symptoms are present, the patient must call for emergency services or proceed to the nearest emergency department immediately. Emergency contact numbers include:
\begin{itemize}
  \item \textbf{local emergency services} in many countries.
  \item \textbf{911} in North America (United States and Canada).
  \item \textbf{999} in the United Kingdom.
  \item \textbf{112} in the European Union and as a global standard on mobile networks.
\end{itemize}

Additionally, patients should seek prompt medical attention if a finger becomes completely locked in a bent position and cannot be straightened with reasonable, gentle passive force, as leaving the joint locked can lead to permanent contractures.

\section*{Diagnosis and Investigations}
The diagnosis of trigger finger is primarily clinical, based on a comprehensive medical history and a targeted physical examination of the hand. In the majority of cases, diagnostic imaging and laboratory investigations are not required to confirm the diagnosis, but they can be invaluable in ruling out concurrent pathologies or assessing underlying systemic diseases.

\begin{itemize}
  \item \textbf{Clinical History}: The physician will inquire about the onset of symptoms, the specific fingers involved, the presence of clicking or locking, and whether the digit requires manual straightening. The clinician will also screen for risk factors, such as diabetes, thyroid disease, inflammatory arthritis, or occupational activities that require repetitive gripping.
  \item \textbf{Physical Examination}:
  \begin{itemize}
    \item \textit{Inspection}: The hand is inspected for signs of swelling, muscular atrophy, or abnormal resting postures.
    \item \textit{Palpation}: The clinician palpatively evaluates the palm at the level of the A1 pulley. Tenderness over the A1 pulley is the most consistent clinical sign. The clinician also feels for a discrete, palpable nodule that moves longitudinally during digit movement.
    \item \textit{Provocative Testing}: The patient is asked to repeatedly flex their fingers into a tight fist and then fully extend them. This allows the clinician to directly observe and feel the catching, clicking, or locking of the tendon.
    \item \textit{Range of Motion Assessment}: The passive and active range of motion of the MCP, PIP, and DIP joints are measured. The presence of any fixed flexion deformity at the PIP joint is documented.
  \end{itemize}
  \item \textbf{Imaging Studies}:
  \begin{itemize}
    \item \textit{Plain Radiographs}: Standard anteroposterior (AP) and lateral X-rays of the hand are typically normal in patients with trigger finger. They are primarily used to exclude other causes of pain or mechanical block, such as osteoarthritis of the MCP or PIP joints, intra-articular loose bodies, calcified deposits, or occult fractures.
    \item \textit{High-Resolution Musculoskeletal Ultrasound}: Ultrasound is a dynamic, non-invasive imaging modality that can confirm trigger finger in atypical cases. Key ultrasound findings include thickening of the A1 pulley (typically greater than 1.0 to 1.5 mm), fluid accumulation (halo sign) within the flexor tendon sheath, tendon swelling (tendinosis), and hyperaemia on Doppler imaging. Most importantly, dynamic ultrasound allows real-time visualisation of the tendon nodule bunching up and catching at the entrance of the A1 pulley during active finger extension.
  \end{itemize}
  \item \textbf{Laboratory Studies}: Blood tests are only indicated if the clinician suspects an underlying systemic condition. If indicated, tests may include Haemoglobin A1c (HbA1c) to screen for diabetes mellitus; Rheumatoid Factor (RF), Anti-Cyclic Citrullinated Peptide (anti-CCP), Erythrocyte Sedimentation Rate (ESR), and C-Reactive Protein (CRP) to evaluate for rheumatoid arthritis or other inflammatory arthropathies; and Thyroid-Stimulating Hormone (TSH) to screen for hypothyroidism.
\end{itemize}

\section*{Management}
The management of trigger finger ranges from conservative, non-invasive measures for mild or early-stage cases to surgical intervention for persistent or severe symptoms. The choice of treatment depends on the duration of symptoms, the grade of triggering, patient preferences, and the presence of underlying systemic conditions.

\begin{itemize}
  \item \textbf{Conservative and Non-Operative Treatments}:
  \begin{itemize}
    \item \textit{Activity Modification}: The first step in management is to identify and modify activities that exacerbate the symptoms. Patients should avoid repetitive gripping, pinching, or using tools that transmit vibration. Ergonomic adjustments, such as using larger handles or padded gloves, can help reduce stress on the flexor tendons.
    \item \textit{Orthotic Splinting}: Splinting is designed to prevent the tendon nodule from gliding through the A1 pulley, thereby reducing friction and allowing the tissue to heal. The metacarpophalangeal (MCP) joint is typically splinted in extension (0 to 15 degrees of flexion) while leaving the PIP and DIP joints free to prevent stiffness. Splinting is worn continuously or at night for a duration of 6 to 10 weeks. Success rates range from 50\% to 70\%, with the highest success observed in patients with short-duration, mild symptoms.
    \item \textit{Pharmacotherapy}: Oral non-steroidal anti-inflammatory drugs (NSAIDs) such as ibuprofen, naproxen, or selective COX-2 inhibitors (e.g., celecoxib) can help manage pain and mild swelling. However, they are generally not curative, as the underlying pathology is degenerative rather than purely inflammatory.
    \item \textit{Corticosteroid Injections}: An injection of a corticosteroid (e.g., triamcinolone acetonide or methylprednisolone) mixed with a local anaesthetic (e.g., lidocaine) directly into the flexor tendon sheath at the level of the A1 pulley. This is highly effective, offering temporary or permanent resolution in 60\% to 90\% of cases. The success rate is lower in patients with long-standing symptoms, multiple affected digits, or diabetes mellitus. Typically, a maximum of two injections per digit is recommended before proceeding to surgery, as repeated injections can weaken the tendon and increase the risk of tendon rupture.
  \end{itemize}
  \item \textbf{Surgical Treatments}:
  \begin{itemize}
    \item \textit{Open Trigger Finger Release (Open TFR)}: Open TFR is the gold-standard surgical treatment, performed as an outpatient procedure under local anaesthesia (often using the Wide Awake Hand Surgery Under Local Anesthesia Only, or WALANT, technique). A small (1 to 2 cm) longitudinal or transverse incision is made in the palm over the A1 pulley. The surgeon carefully dissects the subcutaneous tissue and palmar fascia, taking care to identify and protect the adjacent digital nerves and vessels. The A1 pulley is identified and divided longitudinally. The patient is asked to actively move the finger to confirm that all catching is resolved and that the tendon glides smoothly. The wound is irrigated and closed with sutures.
    \item \textit{Percutaneous Trigger Finger Release}: In this technique, a needle or a specialised micro-scalpel is inserted through the skin directly into the A1 pulley. Using anatomical landmarks, palpation, or real-time ultrasound guidance, the surgeon sweeps the needle or blade longitudinally to divide the A1 pulley. This technique eliminates the need for a surgical incision, resulting in minimal post-operative pain and a faster return to daily activities. However, because the procedure is performed without direct visualisation, it carries a slightly higher risk of incomplete release or injury to the digital nerves, especially in the thumb and index finger where the digital nerves run very close to the pulley.
  \end{itemize}
  \item \textbf{Post-operative Recovery and Hand Therapy}:
  Patients are encouraged to perform gentle, active range-of-motion exercises immediately after surgery to prevent the formation of adhesions between the tendon and the surrounding tissues. The surgical dressing is kept clean and dry, and sutures are removed in 10 to 14 days. While most patients regain full hand function rapidly, hand therapy may be recommended for individuals who develop persistent stiffness, have a pre-existing PIP joint contracture, or have multiple released digits.
\end{itemize}

\section*{Complications}
While trigger finger release is a highly successful and routinely performed procedure, both surgical interventions and leaving the condition untreated carry inherent risks. Understanding these potential complications is essential for informed patient choice and optimal clinical management.

\begin{itemize}
  \item \textbf{Surgical and Procedural Complications}:
  \begin{itemize}
    \item \textit{Incomplete Release}: This is the most common cause of surgical failure. If the A1 pulley is not completely divided, or if the surgeon fails to recognise a co-existing mechanical block (such as thickening of the proximal portion of the A2 pulley or a slip of the FDS tendon), the patient will experience persistent catching or locking after surgery. This typically requires revision surgery.
    \item \textit{Digital Nerve Injury}: The digital nerves run immediately adjacent to the flexor tendon sheath. During open or percutaneous release, these nerves can be stretched, bruised, or accidentally transected. A nerve injury can result in numbness, paresthesia, or burning pain along the side of the finger, or lead to the formation of a painful neuroma, which is challenging to treat.
    \item \textit{Bowstringing}: If the surgeon extends the release too far distally and accidentally damages or cuts the A2 pulley, the flexor tendons will pull away from the bone (bowstringing) during finger flexion. Bowstringing alters the mechanics of the hand, leading to a permanent reduction in active flexion range of motion and a significant loss of grip strength.
    \item \textit{Infection}: Superficial wound infections or deep space hand infections can occur post-operatively. These present with redness, swelling, warmth, and purulent drainage, and require prompt treatment with antibiotics or surgical irrigation.
    \item \textit{Stiffness and Adhesions}: If the patient does not perform early range-of-motion exercises, scar tissue can form between the flexor tendons and the surrounding sheath. This can lead to tendon adhesions, resulting in persistent stiffness, pain, and restricted movement.
    \item \textit{Scar Tenderness and Hypersensitivity}: The surgical scar in the palm is prone to becoming firm, tender, and hypersensitive during the first few weeks to months after surgery. This is generally self-limiting but can be managed with scar massage and desensitisation techniques.
    \item \textit{Complex Regional Pain Syndrome (CRPS)}: A rare but severe complication characterised by chronic, burning pain, swelling, temperature changes, and stiffness in the hand that is disproportionate to the surgical trauma. CRPS requires prompt, multidisciplinary management, including physical therapy, nerve blocks, and pharmacological treatment.
  \end{itemize}
  \item \textbf{Complications of Untreated Trigger Finger}:
  \begin{itemize}
    \item \textit{Fixed Flexion Contractures}: If a finger remains chronically locked or bent, the joints (especially the PIP joint) will lose their normal mobility. Over time, the volar plate, collateral ligaments, and joint capsule shorten and become fibrotic, resulting in a permanent flexion contracture that cannot be corrected even after a delayed tendon release.
    \item \textit{Tendon Wear, Fraying, and Rupture}: The chronic friction and mechanical pressure of the tendon passing through the narrow pulley can cause progressive wear, attenuation, and eventual rupture of the FDS or FDP tendons.
    \item \textit{Functional Atrophy and Disability}: Chronic pain and locking lead to disuse of the hand, which can cause muscle atrophy in the forearm and hand, significantly impairing the patient's ability to perform activities of daily living.
  \end{itemize}
\end{itemize}

\section*{Prognosis and Outcomes}
The overall prognosis for patients with trigger finger is excellent, with high rates of successful treatment and low rates of recurrence.

For patients undergoing non-operative management, conservative treatments can be curative. Corticosteroid injections have a success rate of 60\% to 90\% for a single digit, although recurrence rates are higher in patients with diabetes, multiple affected fingers, or symptoms that have lasted for more than six months. If a second injection is required, the success rate drops, and surgical release is often recommended to prevent tendon damage.

Surgical release, whether open or percutaneous, offers a definitive cure in over 95\% of cases. Recurrence in a successfully released digit is extremely rare.

\begin{itemize}
  \item \textbf{Recovery Timeline}:
  \begin{itemize}
    \item \textit{Immediate}: The mechanical locking and catching of the finger are resolved immediately following the surgical division of the A1 pulley.
    \item \textit{1 to 2 Weeks}: The surgical wound heals, and sutures are removed. Most patients can perform light activities of daily living, such as writing and eating, within a few days of surgery.
    \item \textit{4 to 8 Weeks}: Localised palmar swelling, bruising, and scar tenderness gradually resolve. Patients can begin returning to heavier activities, such as lifting or sports, as tolerated.
    \item \textit{2 to 3 Months}: Hand strength, grip strength, and complete joint flexibility are typically restored to normal.
  \end{itemize}
  \item \textbf{Factors Affecting Outcomes}:
  \begin{itemize}
    \item \textit{Diabetes Mellitus}: Diabetic patients tend to have a higher rate of postoperative stiffness, slower wound healing, and a higher risk of developing trigger finger in other digits.
    \item \textit{Duration of Preoperative Symptoms}: Patients who have had a chronically locked finger for months before undergoing surgery are at a higher risk of persistent stiffness due to joint capsule contractures, which may require prolonged postoperative hand therapy.
    \item \textit{Adherence to Rehabilitation}: Compliance with early range-of-motion exercises is a critical factor in avoiding tendon adhesions and achieving full range of motion.
  \end{itemize}
\end{itemize}

\section*{Prevention and Self-Care}
While it is not always possible to prevent trigger finger, particularly when it is associated with systemic diseases, implementing appropriate self-care measures and ergonomic modifications can significantly reduce the risk of onset, prevent recurrence, and manage early symptoms.

\begin{itemize}
  \item \textbf{Ergonomic Adjustments}:
  \begin{itemize}
    \item Modify tools to reduce the force required to grip them. For example, add foam padding to tool handles, pens, and steering wheels.
    \item Ensure that workstations are designed to maintain a neutral wrist and hand position, avoiding extreme flexion or extension during tasks like typing or assembly line work.
    \item Use tools with a trigger mechanism that distributes the pressure across multiple fingers rather than concentrating the force on a single digit.
  \end{itemize}
  \item \textbf{Activity Management}:
  \begin{itemize}
    \item Avoid prolonged, repetitive hand movements that involve forceful gripping or pinching.
    \item Integrate frequent rest breaks during activities such as gardening, crafting, playing musical instruments, or using hand tools.
    \item Alternate tasks throughout the day to vary the physical demands on the hands.
  \end{itemize}
  \item \textbf{Hand Stretching and Exercises}:
  \begin{itemize}
    \item Perform regular hand stretches to maintain flexibility in the tendons and joints. A simple stretch involves placing the hand flat on a table and slowly lifting each finger individually, or using the opposite hand to gently push the fingers back into extension.
    \item Perform tendon gliding exercises: flex the fingers to form a hook fist, then a straight fist, and finally a full fist, repeating this sequence several times to promote smooth tendon excursion.
  \end{itemize}
  \item \textbf{Home Care for Early Symptoms}:
  \begin{itemize}
    \item Apply ice packs wrapped in a thin towel to the palm for 10 to 15 minutes, three to four times a day, to reduce localised swelling and discomfort when symptoms first appear.
    \item Use warm water soaks or a heating pad in the morning to relax stiff tendons and improve mobility before starting daily activities.
    \item Rest the hand and consider wearing a temporary, over-the-counter finger splint at night to prevent the finger from curling into a fist during sleep.
  \end{itemize}
  \item \textbf{Systemic Health Management}:
  Manage predisposing health conditions effectively. Patients with diabetes should strive to maintain optimal glycaemic control, and those with rheumatoid arthritis should adhere to their prescribed medical regimens to minimise systemic inflammation that can affect the tendon sheaths.
\end{itemize}

\section*{Sources and Further Reading}
For more information, clinical guidelines, and patient support resources, consult the following authoritative organisations and reference guides:

\begin{itemize}
  \item \textbf{American Academy of Orthopaedic Surgeons (AAOS)}: Provides a comprehensive patient education resource outlining the symptoms, causes, diagnosis, and non-surgical and surgical treatments of trigger finger. Available at: \texttt{https://orthoinfo.aaos.org/en/diseases--conditions/trigger-finger}
  \item \textbf{American Society for Surgery of the Hand (ASSH)}: Offers detailed clinical information, diagrams of hand anatomy, and an overview of trigger finger release procedures written by hand surgeons. Available at: \texttt{https://www.assh.org/handcare/condition/trigger-finger}
  \item \textbf{Mayo Clinic}: A trusted clinical overview detailing the diagnostic criteria, risk factors, and therapeutic options for stenosing tenosynovitis. Available at: \texttt{https://www.mayoclinic.org/diseases-conditions/trigger-finger/symptoms-causes/syc-20365100}
  \item \textbf{National Health Service (NHS), United Kingdom}: Public health guide on trigger finger symptoms, non-surgical options, steroid injections, and surgical recovery expectations. Available at: \texttt{https://www.nhs.uk/conditions/trigger-finger/}
  \item \textbf{Healthdirect Australia}: A government-backed health information service providing guidance on trigger finger symptoms, treatment pathways, and finding hand therapy specialists. Available at: \texttt{https://www.healthdirect.gov.au/trigger-finger}
  \item \textbf{The Journal of Hand Surgery}: A leading peer-reviewed journal publishing clinical studies, surgical outcomes, and technical guide updates on open and percutaneous trigger finger release. Available at: \texttt{https://www.jhandsurg.org}
\end{itemize}